% ============================================================
% HÌNH VẼ CHƯƠNG 3 (TikZ) — gọi bằng macro trong Chap3.tex
%   \HinhUseCase     : biểu đồ ca sử dụng tổng quát        (fig:usecase)
%   \HinhOnOffChain  : phân tách dữ liệu on/off-chain       (fig:onoffchain)
%   \HinhTrangThai   : biểu đồ trạng thái của lô tín chỉ    (fig:trangthai)
% Đây là bản đề xuất theo phiếu giao nhiệm vụ: sửa lại cho khớp thiết kế thực tế của bạn.
% ============================================================

% Hình người (tác nhân UML)
\tikzset{
  tacnhan/.pic={
    \draw[thick] (0,0.62) circle (0.15);
    \draw[thick] (0,0.47) -- (0,0.05);
    \draw[thick] (-0.25,0.33) -- (0.25,0.33);
    \draw[thick] (0,0.05) -- (-0.18,-0.28);
    \draw[thick] (0,0.05) -- (0.18,-0.28);
  }
}

% ---------- Biểu đồ ca sử dụng tổng quát ----------
\newcommand{\HinhUseCase}{%
\begin{figure}[htbp]
\centering
\begin{tikzpicture}[font=\footnotesize]
  % Khung hệ thống
  \draw[thick] (3.0,-4.3) rectangle (9.8,4.6);
  \node[font=\footnotesize\bfseries] at (6.4,4.25) {Hệ thống quản lý tín chỉ carbon};
  % Ca sử dụng: tên nút / tung độ / nội dung
  \foreach \ten/\y/\noidung in {
      uc1/3.5/{Thẩm định, phê duyệt và mint tín chỉ},
      uc2/2.6/{Đăng ký dự án},
      uc3/1.7/{Gửi yêu cầu cấp tín chỉ},
      uc4/0.8/{Nạp tín chỉ vào pool, đổi token},
      uc5/-0.1/{Mua bán, chuyển nhượng tín chỉ},
      uc6/-1.0/{Hủy tín chỉ để bù trừ (retire)},
      uc7/-1.9/{Tra cứu tín chỉ, chứng nhận hủy},
      uc8/-2.8/{Theo dõi dashboard hạn ngạch},
      uc9/-3.7/{Quản lý tài khoản, phân quyền}} {
    \draw (6.4,\y) ellipse (3.2cm and 0.38cm);
    \node at (6.4,\y) {\noidung};
    \coordinate (\ten-w) at (3.2,\y);
    \coordinate (\ten-e) at (9.6,\y);
  }
  % Tác nhân bên trái (tên ở dưới)
  \pic at (1.2,3.1) {tacnhan};   \node[align=center] at (1.2,2.55) {VVB};
  \pic at (1.2,0.8) {tacnhan};   \node[align=center] at (1.2,0.25) {Chủ dự án};
  % Tác nhân bên phải (tên ở bên phải để không đè lên đường nối)
  \pic at (11.2,0.5) {tacnhan};  \node[align=left, anchor=west] at (11.5,0.55) {Doanh nghiệp\\mua};
  \pic at (11.2,-3.0) {tacnhan}; \node[align=left, anchor=west] at (11.5,-2.95) {Cơ quan\\quản lý};
  % Liên kết
  \draw (1.5,3.4) -- (uc1-w);
  \foreach \u in {uc2,uc3,uc4,uc5} \draw (1.5,1.1) -- (\u-w);
  \foreach \u in {uc4,uc5,uc6,uc7} \draw (10.9,0.8) -- (\u-e);
  \foreach \u in {uc7,uc8,uc9} \draw (10.9,-2.7) -- (\u-e);
\end{tikzpicture}
\caption{Biểu đồ ca sử dụng tổng quát của hệ thống}
\label{fig:usecase}
\end{figure}}

% ---------- Phân tách dữ liệu on-chain / off-chain ----------
\newcommand{\HinhOnOffChain}{%
\begin{figure}[htbp]
\centering
\begin{tikzpicture}[
    >=Stealth,
    o/.style={draw, rounded corners=2pt, align=left, font=\footnotesize,
              text width=4.7cm, inner sep=4pt, fill=white},
    nhan/.style={font=\scriptsize, align=center, text width=2.6cm}]
  % Off-chain
  \node[o] (files) at (2.6,1.6) {\textbf{Kho tệp}\\
      IPFS hoặc lưu trữ đối tượng:\\
      PDD, báo cáo thẩm định,\\ tệp metadata JSON};
  \node[o, below=5mm of files] (db) {\textbf{CSDL quan hệ}\\
      PostgreSQL, theo mô hình\\ CAD Trust V2.0:\\
      Project, Stakeholders,\\ Verifications, Issuances, Units;\\
      tài khoản ứng dụng (bảng riêng);\\
      lịch sử giao dịch đồng bộ từ sự kiện};
  \node[draw, thick, dashed, rounded corners=4pt, inner sep=6pt, fit=(files)(db),
        label={[font=\footnotesize\bfseries]above:Off-chain (backend)}] (off) {};
  % On-chain
  \node[o, fill=gray!10] (nft) at (11.0,1.6) {\textbf{Hợp đồng NFT tín chỉ (ERC-721)}\\
      tokenId, chủ sở hữu, số sê-ri,\\ mã băm hồ sơ, URI metadata};
  \node[o, fill=gray!10, below=4mm of nft] (pool) {\textbf{Hợp đồng pool (ERC-20)}\\
      số dư token,\\ NFT đang khóa trong pool};
  \node[o, fill=gray!10, below=4mm of pool] (ret) {\textbf{Module Retirement}\\
      sự kiện Retired, người thụ hưởng};
  \node[draw, thick, rounded corners=4pt, inner sep=6pt, fit=(nft)(pool)(ret),
        label={[font=\footnotesize\bfseries]above:On-chain (Ethereum Sepolia)}] (on) {};
  % Luồng dữ liệu
  \draw[->, thick] (files.east) -- node[nhan, above] {mã băm hồ sơ (keccak256) và URI khi mint}
      (on.west |- files.east);
  \draw[->, thick] (on.west |- db.east) -- node[nhan, below] {sự kiện Minted, Transfer, Retired: backend đồng bộ}
      (db.east);
\end{tikzpicture}
\caption{Phân tách dữ liệu giữa phần on-chain và off-chain của hệ thống}
\label{fig:onoffchain}
\end{figure}}

% ---------- Biểu đồ trạng thái của một lô tín chỉ ----------
\newcommand{\HinhTrangThai}{%
\begin{figure}[htbp]
\centering
\begin{tikzpicture}[
    >=Stealth,
    tt/.style={draw, rounded corners=6pt, align=center, font=\footnotesize,
               minimum height=0.9cm, inner xsep=5pt, fill=white},
    nhan/.style={font=\scriptsize, align=center}]
  \node[circle, fill=black, inner sep=3pt] (bd) at (0,0) {};
  \node[tt] (nhap)   at (1.4,0)   {Nháp};
  \node[tt] (cho)    at (4.5,0)   {Chờ thẩm định};
  \node[tt] (tuchoi) at (4.5,-2.3) {Bị từ chối};
  \node[tt] (dacap)  at (9.0,0)   {Đã cấp\\(NFT ERC-721)};
  \node[tt] (pool)   at (9.0,-2.6) {Trong pool\\(ERC-20)};
  \node[tt] (dahuy)  at (12.6,0)   {Đã hủy\\(retired)};
  \node[circle, draw, inner sep=4pt] (kt) at (14.1,0) {};
  \node[circle, fill=black, inner sep=2pt] at (kt) {};

  \draw[->] (bd) -- (nhap);
  \draw[->] (nhap) -- node[nhan, above] {gửi\\hồ sơ} (cho);
  \draw[->] (cho) -- node[nhan, above] {VVB duyệt,\\mint NFT} (dacap);
  \draw[->] (cho) -- node[nhan, right] {VVB từ chối} (tuchoi);
  \draw[->] (tuchoi.west) -| node[nhan, pos=0.3, below] {bổ sung hồ sơ} (nhap.south);
  \draw[->] (dacap) -- node[nhan, above] {retire\\(burn)} (dahuy);
  \draw[->] (dahuy) -- (kt);
  \draw[->] ([xshift=-5mm]dacap.south) -- node[nhan, left] {nạp vào pool} ([xshift=-5mm]pool.north);
  \draw[->] ([xshift=5mm]pool.north) -- node[nhan, right] {đổi lại\\(redeem)} ([xshift=5mm]dacap.south);
  \draw[->] ([xshift=-4mm]dacap.north) .. controls +(-0.3,0.9) and +(0.3,0.9) ..
      node[nhan, above] {chuyển nhượng} ([xshift=4mm]dacap.north);
\end{tikzpicture}
\caption{Biểu đồ trạng thái của một lô tín chỉ carbon trong hệ thống}
\label{fig:trangthai}
\end{figure}}
